\documentclass[../../main]{subfiles}

\begin{document}

\section{Estruturas Algébricas como Modelos de Assinaturas}
\label{sec:matematica:algebra:estruturas-algebricas:modelos}

Nesta seção formalizamos o conceito de estrutura algébrica como uma interpretação de uma assinatura, juntamente com um conjunto de axiomas que restringem o comportamento das operações.

\subsection{Interpretação de assinaturas}

Seja $\Sigma = \{\sigma_i\}_{i \in I}$ uma assinatura algébrica, onde cada símbolo $\sigma_i$ possui aridade $n_i$.

Uma \textbf{interpretação} de $\Sigma$ em um conjunto não vazio $A$ é uma família de funções
\[
  \sigma_i^A : A^{n_i} \to A,
\]
que associa a cada símbolo de operação uma função concreta em $A$.

Essa interpretação substitui os símbolos abstratos por operações efetivas no conjunto.

\subsection{Estrutura algébrica}

Uma \textbf{estrutura algébrica} é um par
\[
  (A, \Sigma),
\]
onde:

\begin{itemize}
  \item $A$ é um conjunto não vazio;
  \item $\Sigma$ é uma assinatura;
  \item cada símbolo $\sigma \in \Sigma$ possui uma interpretação em $A$.
\end{itemize}

Equivalentemente, pode-se escrever a estrutura como
\[
  (A, (\sigma_i^A)_{i \in I}).
\]

\subsection{Termos algébricos}

A partir de uma assinatura $\Sigma$, é possível construir expressões formais chamadas \textbf{termos algébricos}.

Esses termos são obtidos recursivamente a partir de variáveis e aplicações de símbolos de operação.

Exemplos de termos:
\begin{itemize}
  \item $\sigma(x,y)$;
  \item $\sigma(x, \sigma(y,z))$;
  \item expressões compostas de forma finita.
\end{itemize}

Os termos representam combinações formais de operações sem ainda depender de uma interpretação específica.

\subsection{Axiomas algébricos}

Um \textbf{axioma algébrico} é uma sentença da forma equacional, isto é, uma igualdade entre dois termos construídos a partir da assinatura.

Exemplo geral:
\[
  t_1(x_1,\dots,x_n) = t_2(x_1,\dots,x_n).
\]

Os axiomas impõem restrições sobre as interpretações das operações em um conjunto.

\subsection{Modelos de uma assinatura}

Seja $\Sigma$ uma assinatura e $\mathcal{A}$ um conjunto de axiomas sobre $\Sigma$.

Uma estrutura algébrica $(A, \Sigma)$ é um \textbf{modelo} de $\mathcal{A}$ se todas as igualdades especificadas em $\mathcal{A}$ são satisfeitas em $A$.

Escreve-se:
\[
  (A, \Sigma) \vDash \mathcal{A}.
\]

\subsection{Separação sintática e semântica}

A teoria das estruturas algébricas separa dois níveis fundamentais:

\begin{itemize}
  \item \textbf{nível sintático:} assinatura $\Sigma$ e axiomas $\mathcal{A}$, que descrevem formalmente as regras;
  \item \textbf{nível semântico:} conjunto $A$ e interpretações das operações, que realizam concretamente a estrutura.
\end{itemize}

Essa separação permite comparar estruturas distintas que obedecem aos mesmos axiomas.

\subsection{Observação estrutural}

A noção de estrutura algébrica como modelo de uma assinatura permite unificar diversas construções matemáticas sob um mesmo formalismo.

Grupos, anéis, corpos e outras estruturas serão definidos posteriormente como modelos de assinaturas específicas com conjuntos particulares de axiomas.

\end{document}
